% ============================================================================
% Abstract — institution-neutral; shared between the standalone main report and
% the complete dossier (policy-EN.tex).  Content only; the parent document
% provides the heading (\section*/\chapter*) and surrounding pagination.
% ============================================================================

\noindent
Following a period of observation and analysis of the digital policies related to education, digital technology and security at a
research-intensive university, and after reviewing the latest measures under
consideration, it appeared necessary to produce this document in order to
contribute constructively to the evolution of such an institution's digital policy.

\medskip
\noindent
Universities are unique institutions combining cutting-edge research,
advanced teaching and administration. A cybersecurity policy
could tend to apply uniform and over restrictive measures that can
prove counter-productive and at odds with internationally recognised
best practices. This report presents a state of the art of
infrastructure architectures and IT governance models adopted by
major world universities. It shows that \keyword{network segmentation},
\keyword{federated governance} and a \keyword{risk-based approach}
make it possible to reconcile enhanced security and academic freedom, by
drawing on concrete examples from leading universities
(MIT, Stanford, Cambridge, ETH Zurich, Oxford, CERN, etc.), on
international normative frameworks (NIST, ISO~27001, COBIT), and on the
European regulatory context (GDPR Article~32, NIS2 Directive, AI Act).
The report also addresses the issues related to staff BYOD,
artificial intelligence and access to global cloud services
for research.
